\setchapvignette[6cm]{appendix-cookbook}
\chapter{Command cookbook}
\label{app:cookbook}

There are two interchangeable local toolchains for computing and checking \textsc{swhid}s, and
they produce the \emph{same} identifiers for the same bytes --- pick whichever suits your
environment. \textbf{Path~A} is \texttt{swhid-rs}, the Rust reference implementation of the
standard (ISO/IEC~18670, MIT license): fast, standalone, and the one that tracks the specification
most closely. \textbf{Path~B} is the original Python implementation --- \texttt{swh identify}
(from \texttt{swh.model}) and \texttt{swh scanner} --- which adds two things the Rust CLI does not
yet cover: snapshot identifiers, and scanning a working tree against the archive. Appendix~%
\ref{app:swhid-tools} compares all the implementations in full. The two paths below mirror each
other operation for operation; the shared section at the end is the same whichever you chose.

\paragraph{Path A --- the Rust reference tool (\texttt{swhid-rs}).}
\begin{verbatim}
# --- Install -------------------------------------------------------------
cargo install swhid                  # swhid-rs (Rust reference impl., MIT)
cargo install swhid --features git   # ...with VCS (revision) support
# or download a pre-built binary from the project's releases page:
#   https://github.com/swhid/swhid-rs

# --- Compute core SWHIDs locally ----------------------------------------
swhid content --file file.c          # -> swh:1:cnt:...
swhid dir src/                 # -> swh:1:dir:...  (Merkle hash of tree)
swhid parse "swh:1:cnt:<hex>"        # validate / pretty-print a SWHID

# A revision SWHID is just the git commit hash:
echo "swh:1:rev:$(git rev-parse HEAD)"     # -> swh:1:rev:...

# --- Verify bytes against a claimed SWHID (exit 0 iff they match) --------
swhid verify file.c swh:1:cnt:<hex>        # PATH first, then SWHID
swhid verify src/   swh:1:dir:<hex>

# Snapshot SWHIDs and tree scanning: use Path B below.
\end{verbatim}

\paragraph{Path B --- the Python tools (\texttt{swh.model} + \texttt{swh.scanner}).}
\begin{verbatim}
# --- Install -------------------------------------------------------------
pip install 'swh.model[cli]'         # swh identify
pip install swh.scanner              # swh scanner

# --- Compute SWHIDs locally (auto-detects type, or force with --type) ---
swh identify file.c                  # -> swh:1:cnt:...
swh identify --type directory src/   # -> swh:1:dir:...
swh identify --type directory ./export/    # a CLEAN export -> swh:1:dir:...
#   (never a live checkout: see the caveat below)
swh identify --type snapshot repo.git/     # -> swh:1:snp:...  (snapshot)

# --- Verify bytes against a claimed SWHID (exit 0 iff they match) --------
swh identify --verify swh:1:cnt:<hex> file.c    # SWHID first, then PATH

# --- Scan a working tree against the archive ----------------------------
swh scanner scan ./repo                    # what is already archived?
\end{verbatim}

\paragraph{Shared --- archive and resolve (no install needed).}
\begin{verbatim}
# Archive a public repo:  https://save.softwareheritage.org
# Resolve any SWHID:      https://archive.softwareheritage.org/<swhid>
\end{verbatim}

\noindent\textbf{One caveat for \texttt{dir} hashes.} A directory \textsc{swhid} is computed over
everything below the directory you point the tool at --- including a \texttt{.git/} directory,
build output and editor droppings --- and it depends on file \emph{permissions} as well as
content. To match what the archive stored from a git repository, never hash a live checkout:
compute over a clean export (\texttt{git archive} unpacked into an empty directory), whose hash
equals \verb|git rev-parse <commit>^{tree}| by construction. If a directory check mismatches
where a content check would pass, suspect a polluted working copy first, and permission bits
second.
